\documentclass[11pt,a4paper]{article}

\usepackage[a4paper,margin=2.5cm]{geometry}
\usepackage[utf8]{inputenc}
\usepackage[T1]{fontenc}
\usepackage{lmodern}
\usepackage{titlesec}
\usepackage{fancyhdr}
\usepackage{enumitem}
\usepackage{graphicx}

\titleformat{\section}{\normalfont\Large\bfseries}{\thesection}{0.6em}{}
\titlespacing*{\section}{0pt}{1.1em}{0.5em}

\pagestyle{fancy}
\fancyhf{}
\renewcommand{\headrulewidth}{0pt}
\fancyfoot[C]{\small SENTINEL X --- Orange Summer Challenge 2026 --- \thepage}

\setlength{\parindent}{0pt}
\setlength{\parskip}{0.45em}

\begin{document}

\begin{center}
  \includegraphics[height=2.6cm]{/home/rakine/osc/v1_yolo/docs/logo_icon.png}\\[8pt]
  {\Huge\bfseries SENTINEL X}\\[4pt]
  {\large\itshape Voir avant --- Agir à temps}\\[4pt]
  {\small Tech Impact --- Orange Summer Challenge 2026}
\end{center}

\vspace{1em}

SENTINEL X donne aux caméras de surveillance la capacité de comprendre ce qu'elles filment, et non plus seulement de détecter un mouvement. Au lieu du détecteur de mouvement traditionnel, qui signale une agitation sans savoir ce qu'elle signifie, le système observe en temps réel et reconnaît des situations : une présence prolongée et suspecte, un comportement qui précède un vol, un cambriolage ou une agression. Il donne l'alerte avant que l'acte ne se produise, sur la base d'une véritable analyse du danger.

\section{Le problème}
Les caméras de surveillance actuelles sont passives. Elles enregistrent, mais elles n'analysent rien de ce qui se passe réellement devant elles. Le détecteur de mouvement classique ne fait pas la différence entre un événement anodin et une situation dangereuse : il réagit à un déplacement, jamais à un comportement. Résultat : la caméra ne sert qu'à constater les faits après coup, une fois le crime, le cambriolage ou l'agression déjà commis.

\section{La mission}
Faire passer la vidéosurveillance d'un rôle de témoin, qui ne sert qu'à prouver qu'un acte a eu lieu, à un rôle de vigie, capable d'aider à l'empêcher. Quatre principes guident cette approche :

\begin{enumerate}[leftmargin=1.4em, itemsep=2pt]
  \item \textbf{Observer le comportement, pas l'identité.} Le système s'intéresse à la posture, aux déplacements et aux gestes des personnes filmées.
  \item \textbf{Reconnaître les signaux qui précèdent un acte.} Repérer les comportements caractéristiques d'une situation à risque avant qu'elle ne dégénère.
  \item \textbf{Alerter un humain, ne jamais décider seul.} Le système signale une situation suspecte à un opérateur, qui reste seul décisionnaire.
  \item \textbf{Adapter la réponse à la gravité.} Une simple notification pour un signal faible ; un appel automatisé pour une situation critique.
\end{enumerate}

\section{Où en est le projet}
Un prototype fonctionnel (MVP) démontre déjà cette approche sur un réseau multi-caméra, multi-bâtiment :

\begin{itemize}[leftmargin=1.4em, itemsep=2pt]
  \item Une chaîne de modèles d'analyse comportementale (menace, feu, bagarre, chute, équipement de sécurité, vêtements).
  \item Une reconnaissance faciale avec gestion des autorisations de présence par lieu.
  \item Une fiche par personne (genre, expression, signalement) et une détection de présence prolongée (rôdeur potentiel).
  \item Un assistant conversationnel dont les réponses reposent strictement sur les données réellement observées, jamais une supposition.
\end{itemize}

Reste à construire pour rejoindre l'ambition complète du projet : une détection plus fine des signaux précurseurs d'un acte (pas seulement de sa conséquence), l'escalade automatisée par appel vocal, et un déploiement pilote sur un site réel.

\vspace{1em}
\noindent\textit{SENTINEL X --- un système qui regarde avant que le pire n'arrive.}

\end{document}
